
\documentclass[12pt]{article}

% ---------- Packages ----------
\usepackage{amsmath}
\usepackage{amssymb}
\usepackage{amsthm}
\usepackage{geometry}
\usepackage{enumitem}
\usepackage{hyperref}

\geometry{a4paper, margin=1in}
\raggedbottom

% Allow runs of empty headings to break across pages.
\AddToHook{cmd/section/before}{\par\penalty0}
\AddToHook{cmd/subsection/before}{\par\penalty0}

% ---------- Theorem Environments ----------
\newtheorem{definition}{Definition}
\newtheorem{theorem}{Theorem}
\newtheorem{lemma}{Lemma}
\newtheorem{corollary}{Corollary}
\newtheorem{example}{Example}

% ---------- Document ----------
\begin{document}

\title{\textbf{Numbers, Sets and Functions}}
\author{Annukul Ghale}


\maketitle

\tableofcontents
\newpage


% ==================================================
% WEEK 1
% ==================================================

\section{Week 1}

\subsection{Definitions}
\begin{definition}
Definition 1.1. A proof is a formal logical argument justifying that a mathematical
statement is true, by proceeding from given assumptions to the desired conclusion.
\end{definition}
\begin{definition}
Theorem 1.2. The sum of two even numbers is even.
\end{definition}

\begin{definition}
A statement represents an expression that:
•is a complete sentence by itself, and
•is either true or false.
\end{definition}

\subsection{Notes}

Formal logic can be viewed in two ways. Boolean logic focuses on whether statements
are true or false, which is useful in computing. Deductive logic focuses on how
conclusions follow from assumptions, as in mathematical proofs.

Let $P$, $Q$, and $R$ be statements. The symbols $\land$, $\lor$, and $\neg$
mean ``and,'' ``or,'' and ``not.'' The symbol $\Rightarrow$ means ``implies,''
and $\Leftrightarrow$ means ``is logically equivalent to.''

The following laws give logically equivalent ways to write statements:
\begin{align*}
\text{Commutative laws:}\quad
  P \land Q &\Leftrightarrow Q \land P,\\
  P \lor Q &\Leftrightarrow Q \lor P,\\
\text{Associative laws:}\quad
  P \land (Q \land R) &\Leftrightarrow (P \land Q) \land R,\\
  P \lor (Q \lor R) &\Leftrightarrow (P \lor Q) \lor R,\\
\text{De Morgan's laws:}\quad
  \neg(P \land Q) &\Leftrightarrow (\neg P) \lor (\neg Q),\\
  \neg(P \lor Q) &\Leftrightarrow (\neg P) \land (\neg Q),\\
\text{Distributive laws:}\quad
  P \land (Q \lor R) &\Leftrightarrow (P \land Q) \lor (P \land R),\\
  P \lor (Q \land R) &\Leftrightarrow (P \lor Q) \land (P \lor R).
\end{align*}

An equivalence can be expressed as implications in both directions:
\[
P \Leftrightarrow Q
\quad\text{means}\quad
(P \Rightarrow Q) \land (Q \Rightarrow P).
\]
An implication has the same logical meaning as its contrapositive:
\[
P \Rightarrow Q
\quad\Leftrightarrow\quad
(\neg Q) \Rightarrow (\neg P).
\]




% ==================================================
% WEEK 2
% ==================================================

\section{Week 2}

\subsection{Definitions}
\subsection{Notes}


% ==================================================
% WEEK 3
% ==================================================

\section{Week 3}

\subsection{Definitions}
\subsection{Notes}




% ==================================================
% WEEK 4
% ==================================================

\section{Week 4}

\subsection{Definitions}
\subsection{Notes}




% ==================================================
% WEEK 5
% ==================================================

\section{Week 5}

\subsection{Definitions}
\subsection{Notes}




% ==================================================
% WEEK 6
% ==================================================

\section{Week 6}

\subsection{Definitions}
\subsection{Notes}




% ==================================================
% WEEK 7
% ==================================================

\section{Week 7}

\subsection{Definitions}
\subsection{Notes}



% ==================================================
% WEEK 8
% ==================================================

\section{Week 8}

\subsection{Definitions}
\subsection{Notes}




% ==================================================
% WEEK 9
% ==================================================

\section{Week 9}

\subsection{Definitions}
\subsection{Notes}




% ==================================================
% WEEK 10
% ==================================================

\section{Week 10}

\subsection{Definitions}
\subsection{Notes}




% ==================================================
% WEEK 11
% ==================================================

\section{Week 11}

\subsection{Definitions}
\subsection{Notes}




% ==================================================
% WEEK 12
% ==================================================

\section{Week 12}

\subsection{Definitions}
\subsection{Notes}




\end{document}
